\chapter{PENDAHULUAN}

\section{Latar Belakang}

Dalam administrasi bisnis antarbisnis atau \textit{business-to-business} (B2B), \textit{invoice} merupakan dokumen penagihan yang perlu dikirim, dilacak, dan dibuktikan penerimaannya. Sistem pelacakan mendukung pencatatan status dan riwayat perpindahan dokumen \cite{10.1007/978-3-642-55032-4_33}. Sementara itu, \textit{Proof of Delivery} (POD) menyediakan bukti penerimaan yang diperlukan untuk memverifikasi penyelesaian pengiriman \cite{gunadhya_rajawali_logistik_2021}. Dengan demikian, pelacakan \textit{invoice} dan penyimpanan POD menjadi bagian penting dalam pengendalian administrasi pengiriman dokumen.

Pada perusahaan yang menjadi objek penelitian, pelanggan memiliki jadwal penerimaan dan ketentuan \textit{cutoff} yang berbeda. Kondisi tersebut menuntut admin menentukan \textit{invoice} yang perlu didahulukan. Target keputusan dalam penelitian ini terdiri atas dua kelas, yaitu \textit{Urgent} dan \textit{Not Urgent}. Keputusan prioritas perlu dirumuskan secara konsisten, dapat dijelaskan, dan dapat dievaluasi terhadap keputusan historis admin.

Pendekatan pertama adalah \textit{Rule-Based} yang dibangun dari SOP perusahaan. SOP diperoleh melalui \textit{Knowledge Acquisition}, diterjemahkan menjadi aturan IF--THEN melalui \textit{Rule Construction}, dan diperiksa melalui \textit{Rule Validation}. Setiap prediksi menyimpan Rule ID yang aktif sehingga hubungan antara kondisi operasional, aturan, dan keluaran dapat ditelusuri. Dalam penelitian ini, proses tersebut digunakan untuk membangun dan memvalidasi \textit{Rule-Based}, bukan untuk membentuk label acuan.

Pendekatan kedua adalah \textit{Decision Tree} berbasis \textit{entropy}. Model dilatih menggunakan atribut historis dan label admin yang tersimpan pada \texttt{expert\_label}. \textit{Decision Tree} mempelajari hubungan antara atribut dan kelas \textit{Urgent} atau \textit{Not Urgent}, kemudian menghasilkan jalur keputusan yang dapat diperiksa. Pelatihan dan pemilihan parameter hanya dilakukan pada data latih agar evaluasi menggunakan data yang tidak memengaruhi pembentukan model.

Kumpulan penelitian yang ditinjau membahas pelacakan dokumen \cite{10.1007/978-3-642-55032-4_33}, POD \cite{gunadhya_rajawali_logistik_2021}, prediksi pembayaran \textit{invoice} \cite{firmansyah,dyahikbal}, serta klasifikasi dan prioritas pada konteks sistem pendukung keputusan \cite{Diwandari_Sela_Syahputra_Parama_Septiarini_2023}. Namun, kumpulan penelitian tersebut belum memberikan bukti komparatif langsung mengenai performa, pola kesalahan, dan karakteristik operasional \textit{Rule-Based} berbasis SOP dan \textit{Decision Tree} berbasis data historis admin untuk klasifikasi prioritas pengiriman \textit{invoice}. Perbandingan tersebut memerlukan \textit{ground truth} yang independen dari kedua metode agar evaluasi tidak dibentuk oleh keluaran salah satu metode.

Penelitian ini membangun \textit{Rule-Based} dari SOP dan melatih \textit{Decision Tree} menggunakan data historis berlabel admin. Kedua metode dievaluasi terhadap label historis admin sebagai \textit{ground truth} dan pada kasus validasi yang sama melalui \textit{stratified hold-out} 80:20, \textit{stratified hold-out} 70:30, \textit{stratified 5-fold cross-validation}, dan \textit{leave-one-out cross-validation} (LOOCV). Perbandingan mencakup performa klasifikasi, pola kesalahan, uji statistik berpasangan, dan karakteristik operasional. Hasil klasifikasi ditempatkan sebagai rekomendasi prioritas dalam konteks sistem pelacakan \textit{invoice} dan POD.

\section{Rumusan Masalah}

Berdasarkan latar belakang yang telah diuraikan, rumusan masalah dalam penelitian ini adalah sebagai berikut.

\begin{enumerate}
    \item Bagaimana membangun klasifikasi \textit{Rule-Based} dari SOP perusahaan dan \textit{Decision Tree} berbasis \textit{entropy} dari data historis berlabel admin pada target biner yang sama?
    \item Bagaimana perbandingan performa dan pola kesalahan kedua metode pada skenario E1--E4 ketika dievaluasi terhadap label historis admin yang sama?
    \item Bagaimana perbedaan karakteristik \textit{Rule-Based} dan \textit{Decision Tree} ditinjau dari \textit{explainability}, \textit{maintainability}, \textit{adaptability}, \textit{scalability}, biaya komputasi, transparansi aturan, kemudahan pembaruan, dan kepatuhan terhadap SOP?
\end{enumerate}

\section{Tujuan Penelitian}

Berdasarkan rumusan masalah tersebut, tujuan penelitian ini adalah sebagai berikut.

\begin{enumerate}
    \item Membangun dan memvalidasi \textit{Rule-Based} melalui \textit{Knowledge Acquisition}, \textit{Rule Construction}, dan \textit{Rule Validation}, serta melatih \textit{Decision Tree} berbasis \textit{entropy} menggunakan data historis berlabel admin.
    \item Mengevaluasi dan membandingkan performa serta pola kesalahan kedua metode pada skenario E1--E4 terhadap \textit{ground truth} admin yang sama.
    \item Menganalisis perbedaan karakteristik operasional kedua metode menggunakan indikator \textit{explainability}, \textit{maintainability}, \textit{adaptability}, \textit{scalability}, biaya komputasi, transparansi aturan, kemudahan pembaruan, dan kepatuhan terhadap SOP.
\end{enumerate}

\section{Batasan Masalah}

Agar penelitian memiliki ruang lingkup yang jelas dan terukur, batasan masalah dalam penelitian ini ditetapkan sebagai berikut.

\begin{enumerate}
    \item Unit analisis adalah 99 \textit{invoice} unik dari 101 baris sumber setelah dua duplikat ditangani.
    \item Dataset bersih terdiri atas 30 \textit{invoice Urgent} dan 69 \textit{invoice Not Urgent}.
    \item \textit{Ground truth} berasal dari \texttt{expert\_label} sebagai label historis admin perusahaan.
    \item Penelitian dibatasi pada satu perusahaan dan periode data yang tersedia.
    \item Metode yang dibandingkan adalah \textit{Rule-Based} dan \textit{Decision Tree} berbasis \textit{entropy}.
    \item \textit{Rule-Based} menggunakan aturan R1--R8 yang diturunkan dari SOP dan dibekukan sebelum evaluasi.
    \item Fitur prediktor hanya menggunakan informasi yang tersedia pada waktu keputusan prioritas.
    \item Atribut \texttt{sent\_date}, \texttt{expert\_reason}, identifier, identitas pelanggan, dan \texttt{Driver} tidak digunakan sebagai fitur prediktor.
    \item Evaluasi menggunakan E1 \textit{stratified hold-out} 80:20, E2 \textit{stratified hold-out} 70:30, E3 \textit{stratified 5-fold cross-validation}, dan E4 LOOCV.
    \item Evaluasi kuantitatif meliputi \textit{confusion matrix}, \textit{accuracy}, \textit{precision}, \textit{recall}, \textit{F1-score}, \textit{macro F1}, \textit{false positive}, \textit{false negative}, dan \textit{exact McNemar}.
    \item \textit{Adaptability} dan \textit{scalability} dibahas secara prosedural atau konseptual karena tidak diuji melalui eksperimen perubahan SOP atau peningkatan skala data khusus.
    \item Penelitian tidak melakukan validasi temporal, eksternal, atau lintas perusahaan.
    \item Sistem hanya menjadi konteks penerapan rekomendasi prioritas, pelacakan \textit{invoice}, dan penyimpanan POD.
\end{enumerate}

\section{Kontribusi Penelitian}

Kontribusi penelitian ini adalah sebagai berikut.

\begin{enumerate}
    \item Prosedur pembentukan \textit{Rule-Based} yang dapat ditelusuri dari SOP melalui \textit{Knowledge Acquisition}, \textit{Rule Construction}, dan \textit{Rule Validation}.
    \item Evaluasi komparatif berpasangan antara \textit{Rule-Based} dan \textit{Decision Tree} pada \textit{ground truth} admin dan kasus validasi yang sama.
    \item Analisis yang menggabungkan metrik prediksi, \textit{false positive}, \textit{false negative}, kesalahan per \textit{invoice}, uji McNemar, serta karakteristik operasional kedua metode.
\end{enumerate}
